\documentclass[10pt]{extarticle}
\usepackage[a4paper,margin=1.5cm]{geometry}
\usepackage[T1]{fontenc}
\usepackage{amsmath,amssymb}
\usepackage{graphicx}
\usepackage[export]{adjustbox}
\usepackage{multicol}
\usepackage{enumitem}
\usepackage{booktabs,array,tabularx}
\setlength{\columnsep}{0.7cm}
\setlength{\parindent}{0pt}
\setlength{\parskip}{3pt}
\pagestyle{empty}
\makeatletter
% figure registry (figdefs.tex is generated by the extract step)
\newcommand{\qfigdef}[4]{\@namedef{qf@file@#1}{#2}\@namedef{qf@w@#1}{#3}\@namedef{qf@wide@#1}{#4}}
\newcommand{\qf@check}[1]{\@ifundefined{qf@file@#1}{\PackageError{score2}{Unknown figure id #1}{}}{}}
\newcommand{\qf@img}[4][]{\includegraphics[width=#2\dimexpr\@nameuse{qf@w@#4}pt\relax,max width=#3,max height=0.45\textheight,#1]{\@nameuse{qf@file@#4}}}
% Full-width questions: finalize wraps every question that contains a wide figure in \qfullbegin...\qfullend,
% so the WHOLE question (stem, figures, options) is set across the page in reading order, then the two columns
% resume. The question's head - its number, stem and first figure (\qfullsplit marks the end of the head) -
% is kept on one page, so a number is never left at a page foot with its figure overleaf; the rest may break
% across pages normally, so no big gaps are left. Without a \qfullsplit, the whole question is kept together
% if it is under 45% of a page. Never write these yourself.
\newif\ifqf@full
\newif\ifqf@boxed
\newsavebox{\qf@box}
\newcommand{\qf@open}{\global\qf@fulltrue\par\global\qf@boxedtrue
  \setbox\qf@box\vbox\bgroup\hsize\textwidth\linewidth\textwidth}
\newcommand{\qf@emit}[1]{\par\egroup\global\qf@boxedfalse
  \ifdim\dimexpr\ht\qf@box+\dp\qf@box\relax<#1\textheight\noindent\usebox{\qf@box}\par\else\unvbox\qf@box\fi}
\newcommand{\qfullsplit}{\ifqf@boxed\qf@emit{0.85}\fi}
\newcommand{\qf@close}{\ifqf@boxed\qf@emit{0.45}\fi\global\qf@fullfalse}
\newcommand{\qfullbegin}{\end{multicols}\qf@open}
\newcommand{\qfullend}{\qf@close\begin{multicols}{2}}
\newcommand{\qfullstart}{\qf@open}
\newcommand{\qfulldone}{\qf@close}
% \qfig[scale]{id}: block figure, centred, directly where it is written (size set by the tool).
% Fallback if a wide figure is not inside a full-width question: break out of the columns in place - no floats.
\newcommand{\qfig}[2][1]{\qf@check{#2}%
  \ifqf@full\qf@blk{#1}{#2}\else
  \ifnum\@nameuse{qf@wide@#2}=1 \expandafter\@firstoftwo\else\expandafter\@secondoftwo\fi
  {\end{multicols}\begin{center}\qf@img{#1}{\textwidth}{#2}\end{center}\begin{multicols}{2}}%
  {\qf@blk{#1}{#2}}\fi}
\newcommand{\qf@blk}[2]{\par\smallskip{\centering\qf@img{#1}{\linewidth}{#2}\par}\smallskip}
% \qfigi[scale]{id}: inline figure (inside option lists / tables / minipages); never breaks columns.
\newcommand{\qfigi}[2][1]{\qf@check{#2}\qf@img[valign=c]{#1}{\linewidth}{#2}}
\makeatother
% structure macros - lint reads these to check completeness and figure placement
\newcommand{\papersec}[1]{\par\bigskip{\centering\Large\bfseries PAPER-#1\par}\medskip}
\newcommand{\partsec}[2]{\par\medskip\noindent{\large\bfseries PART-#1: #2}\par\smallskip}
\newcommand{\sect}[1]{\par\smallskip\noindent\textbf{#1}\par}
\newcommand{\qn}[1]{\par\smallskip\noindent\textbf{#1.}~}
\newcommand{\paraq}[2]{\par\smallskip\noindent\textbf{Paragraph for Questions #1 and #2}\par}
\newenvironment{qoptions}{\begin{itemize}[leftmargin=1.8em,itemsep=1pt,topsep=2pt,label={}]}{\end{itemize}}
\newcommand{\opt}[2]{\item[(#1)] #2}
% figure registry (generated by score2_tools.py finalize)
\qfigdef{p1_phy_s1_q01_1}{Enthusiast Pre-Score - Test 03 (23 Nov 2020)-img/p1_phy_s1_q01_1.png}{212.9}{0}
\qfigdef{p1_phy_s1_q02_2}{Enthusiast Pre-Score - Test 03 (23 Nov 2020)-img/p1_phy_s1_q02_2.png}{141.8}{0}
\qfigdef{p1_phy_s1_q02_1}{Enthusiast Pre-Score - Test 03 (23 Nov 2020)-img/p1_phy_s1_q02_1.png}{164.1}{0}
\qfigdef{p1_phy_s1_q02_3}{Enthusiast Pre-Score - Test 03 (23 Nov 2020)-img/p1_phy_s1_q02_3.png}{176.9}{0}
\qfigdef{p1_phy_s1_q02_4}{Enthusiast Pre-Score - Test 03 (23 Nov 2020)-img/p1_phy_s1_q02_4.png}{171.9}{0}
\qfigdef{p1_phy_s1_q02_5}{Enthusiast Pre-Score - Test 03 (23 Nov 2020)-img/p1_phy_s1_q02_5.png}{173.5}{0}
\qfigdef{p1_phy_s1_q05_1}{Enthusiast Pre-Score - Test 03 (23 Nov 2020)-img/p1_phy_s1_q05_1.png}{245.2}{0}
\qfigdef{p1_phy_s2_para7_8_1}{Enthusiast Pre-Score - Test 03 (23 Nov 2020)-img/p1_phy_s2_para7_8_1.png}{136.8}{0}
\qfigdef{p1_phy_s2_q07_A}{Enthusiast Pre-Score - Test 03 (23 Nov 2020)-img/p1_phy_s2_q07_A.png}{90.1}{0}
\qfigdef{p1_phy_s2_q07_B}{Enthusiast Pre-Score - Test 03 (23 Nov 2020)-img/p1_phy_s2_q07_B.png}{95.4}{0}
\qfigdef{p1_phy_s2_q07_C}{Enthusiast Pre-Score - Test 03 (23 Nov 2020)-img/p1_phy_s2_q07_C.png}{90.1}{0}
\qfigdef{p1_phy_s2_q07_D}{Enthusiast Pre-Score - Test 03 (23 Nov 2020)-img/p1_phy_s2_q07_D.png}{90.1}{0}
\qfigdef{p1_phy_s2_q08_A}{Enthusiast Pre-Score - Test 03 (23 Nov 2020)-img/p1_phy_s2_q08_A.png}{65.6}{0}
\qfigdef{p1_phy_s2_q08_B}{Enthusiast Pre-Score - Test 03 (23 Nov 2020)-img/p1_phy_s2_q08_B.png}{68.9}{0}
\qfigdef{p1_phy_s2_q08_C}{Enthusiast Pre-Score - Test 03 (23 Nov 2020)-img/p1_phy_s2_q08_C.png}{86.4}{0}
\qfigdef{p1_phy_s2_q08_D}{Enthusiast Pre-Score - Test 03 (23 Nov 2020)-img/p1_phy_s2_q08_D.png}{86.4}{0}
\qfigdef{p1_phy_s2_para9_10_1}{Enthusiast Pre-Score - Test 03 (23 Nov 2020)-img/p1_phy_s2_para9_10_1.png}{103.2}{0}
\qfigdef{p1_phy_s3_q04_1}{Enthusiast Pre-Score - Test 03 (23 Nov 2020)-img/p1_phy_s3_q04_1.png}{169.7}{0}
\qfigdef{p1_phy_s3_q05_1}{Enthusiast Pre-Score - Test 03 (23 Nov 2020)-img/p1_phy_s3_q05_1.png}{160.8}{0}
\qfigdef{p1_phy_s3_q06_1}{Enthusiast Pre-Score - Test 03 (23 Nov 2020)-img/p1_phy_s3_q06_1.png}{129.6}{0}
\qfigdef{p1_phy_s3_q07_1}{Enthusiast Pre-Score - Test 03 (23 Nov 2020)-img/p1_phy_s3_q07_1.png}{60.2}{0}
\qfigdef{p1_phy_s3_q08_1}{Enthusiast Pre-Score - Test 03 (23 Nov 2020)-img/p1_phy_s3_q08_1.png}{98.2}{0}
\qfigdef{p1_chem_s1_q02_1}{Enthusiast Pre-Score - Test 03 (23 Nov 2020)-img/p1_chem_s1_q02_1.png}{292.5}{1}
\qfigdef{p1_chem_s1_q05_A}{Enthusiast Pre-Score - Test 03 (23 Nov 2020)-img/p1_chem_s1_q05_A.png}{96.9}{0}
\qfigdef{p1_chem_s1_q05_B}{Enthusiast Pre-Score - Test 03 (23 Nov 2020)-img/p1_chem_s1_q05_B.png}{142.0}{0}
\qfigdef{p1_chem_s1_q05_C}{Enthusiast Pre-Score - Test 03 (23 Nov 2020)-img/p1_chem_s1_q05_C.png}{120.9}{0}
\qfigdef{p1_chem_s1_q05_D}{Enthusiast Pre-Score - Test 03 (23 Nov 2020)-img/p1_chem_s1_q05_D.png}{167.1}{0}
\qfigdef{p1_chem_s1_q06_1}{Enthusiast Pre-Score - Test 03 (23 Nov 2020)-img/p1_chem_s1_q06_1.png}{238.3}{0}
\qfigdef{p1_chem_s2_para7_8_1}{Enthusiast Pre-Score - Test 03 (23 Nov 2020)-img/p1_chem_s2_para7_8_1.png}{209.9}{0}
\qfigdef{p1_chem_s2_q07_A}{Enthusiast Pre-Score - Test 03 (23 Nov 2020)-img/p1_chem_s2_q07_A.png}{65.9}{0}
\qfigdef{p1_chem_s2_q07_B}{Enthusiast Pre-Score - Test 03 (23 Nov 2020)-img/p1_chem_s2_q07_B.png}{54.4}{0}
\qfigdef{p1_chem_s2_q07_C}{Enthusiast Pre-Score - Test 03 (23 Nov 2020)-img/p1_chem_s2_q07_C.png}{54.1}{0}
\qfigdef{p1_chem_s2_q07_D}{Enthusiast Pre-Score - Test 03 (23 Nov 2020)-img/p1_chem_s2_q07_D.png}{41.9}{0}
\qfigdef{p1_chem_s2_q08_A}{Enthusiast Pre-Score - Test 03 (23 Nov 2020)-img/p1_chem_s2_q08_A.png}{209.5}{0}
\qfigdef{p1_chem_s2_q08_B}{Enthusiast Pre-Score - Test 03 (23 Nov 2020)-img/p1_chem_s2_q08_B.png}{139.8}{0}
\qfigdef{p1_chem_s2_q08_C}{Enthusiast Pre-Score - Test 03 (23 Nov 2020)-img/p1_chem_s2_q08_C.png}{147.7}{0}
\qfigdef{p1_chem_s2_q08_D}{Enthusiast Pre-Score - Test 03 (23 Nov 2020)-img/p1_chem_s2_q08_D.png}{185.0}{0}
\qfigdef{p1_chem_s2_q09_A}{Enthusiast Pre-Score - Test 03 (23 Nov 2020)-img/p1_chem_s2_q09_A.png}{124.4}{0}
\qfigdef{p1_chem_s2_q09_B}{Enthusiast Pre-Score - Test 03 (23 Nov 2020)-img/p1_chem_s2_q09_B.png}{142.5}{0}
\qfigdef{p1_chem_s2_q09_C}{Enthusiast Pre-Score - Test 03 (23 Nov 2020)-img/p1_chem_s2_q09_C.png}{128.1}{0}
\qfigdef{p1_chem_s2_q09_D}{Enthusiast Pre-Score - Test 03 (23 Nov 2020)-img/p1_chem_s2_q09_D.png}{94.7}{0}
\qfigdef{p1_chem_s3_q03_1}{Enthusiast Pre-Score - Test 03 (23 Nov 2020)-img/p1_chem_s3_q03_1.png}{409.7}{1}
\qfigdef{p1_chem_s3_q04_1}{Enthusiast Pre-Score - Test 03 (23 Nov 2020)-img/p1_chem_s3_q04_1.png}{234.7}{0}
\qfigdef{p1_math_s1_q02_1}{Enthusiast Pre-Score - Test 03 (23 Nov 2020)-img/p1_math_s1_q02_1.png}{100.4}{0}
\qfigdef{p2_phy_s1_q01_1}{Enthusiast Pre-Score - Test 03 (23 Nov 2020)-img/p2_phy_s1_q01_1.png}{99.3}{0}
\qfigdef{p2_phy_s1_q02_1}{Enthusiast Pre-Score - Test 03 (23 Nov 2020)-img/p2_phy_s1_q02_1.png}{108.5}{0}
\qfigdef{p2_phy_s1_q03_1}{Enthusiast Pre-Score - Test 03 (23 Nov 2020)-img/p2_phy_s1_q03_1.png}{87.7}{0}
\qfigdef{p2_phy_s1_q04_1}{Enthusiast Pre-Score - Test 03 (23 Nov 2020)-img/p2_phy_s1_q04_1.png}{131.4}{0}
\qfigdef{p2_phy_s1_q05_1}{Enthusiast Pre-Score - Test 03 (23 Nov 2020)-img/p2_phy_s1_q05_1.png}{149.5}{0}
\qfigdef{p2_phy_s1_q06_1}{Enthusiast Pre-Score - Test 03 (23 Nov 2020)-img/p2_phy_s1_q06_1.png}{119.4}{0}
\qfigdef{p2_phy_s2_q07_1}{Enthusiast Pre-Score - Test 03 (23 Nov 2020)-img/p2_phy_s2_q07_1.png}{72.9}{0}
\qfigdef{p2_phy_s2_q07_2}{Enthusiast Pre-Score - Test 03 (23 Nov 2020)-img/p2_phy_s2_q07_2.png}{74.2}{0}
\qfigdef{p2_phy_s2_q07_3}{Enthusiast Pre-Score - Test 03 (23 Nov 2020)-img/p2_phy_s2_q07_3.png}{54.4}{0}
\qfigdef{p2_phy_s2_q07_4}{Enthusiast Pre-Score - Test 03 (23 Nov 2020)-img/p2_phy_s2_q07_4.png}{76.4}{0}
\qfigdef{p2_phy_s2_q08_1}{Enthusiast Pre-Score - Test 03 (23 Nov 2020)-img/p2_phy_s2_q08_1.png}{406.3}{1}
\qfigdef{p2_phy_s2_q09_1}{Enthusiast Pre-Score - Test 03 (23 Nov 2020)-img/p2_phy_s2_q09_1.png}{112.2}{0}
\qfigdef{p2_phy_s2_q10_1}{Enthusiast Pre-Score - Test 03 (23 Nov 2020)-img/p2_phy_s2_q10_1.png}{103.6}{0}
\qfigdef{p2_phy_s2_q10_2}{Enthusiast Pre-Score - Test 03 (23 Nov 2020)-img/p2_phy_s2_q10_2.png}{103.6}{0}
\qfigdef{p2_phy_s2_q10_3}{Enthusiast Pre-Score - Test 03 (23 Nov 2020)-img/p2_phy_s2_q10_3.png}{109.1}{0}
\qfigdef{p2_phy_s2_q10_4}{Enthusiast Pre-Score - Test 03 (23 Nov 2020)-img/p2_phy_s2_q10_4.png}{103.6}{0}
\qfigdef{p2_phy_s3_q02_1}{Enthusiast Pre-Score - Test 03 (23 Nov 2020)-img/p2_phy_s3_q02_1.png}{126.4}{0}
\qfigdef{p2_phy_s3_q03_1}{Enthusiast Pre-Score - Test 03 (23 Nov 2020)-img/p2_phy_s3_q03_1.png}{110.0}{0}
\qfigdef{p2_phy_s3_q05_1}{Enthusiast Pre-Score - Test 03 (23 Nov 2020)-img/p2_phy_s3_q05_1.png}{184.5}{0}
\qfigdef{p2_phy_s3_q07_1}{Enthusiast Pre-Score - Test 03 (23 Nov 2020)-img/p2_phy_s3_q07_1.png}{145.7}{0}
\qfigdef{p2_phy_s3_q08_1}{Enthusiast Pre-Score - Test 03 (23 Nov 2020)-img/p2_phy_s3_q08_1.png}{105.6}{0}
\qfigdef{p2_chem_s1_q03_B}{Enthusiast Pre-Score - Test 03 (23 Nov 2020)-img/p2_chem_s1_q03_B.png}{371.8}{1}
\qfigdef{p2_chem_s1_q04_A}{Enthusiast Pre-Score - Test 03 (23 Nov 2020)-img/p2_chem_s1_q04_A.png}{333.2}{1}
\qfigdef{p2_chem_s3_q02_1}{Enthusiast Pre-Score - Test 03 (23 Nov 2020)-img/p2_chem_s3_q02_1.png}{347.5}{1}
\qfigdef{p2_chem_s3_q03_1}{Enthusiast Pre-Score - Test 03 (23 Nov 2020)-img/p2_chem_s3_q03_1.png}{365.3}{1}
\qfigdef{p2_chem_s3_q08_1}{Enthusiast Pre-Score - Test 03 (23 Nov 2020)-img/p2_chem_s3_q08_1.png}{353.0}{1}
\qfigdef{p2_math_s3_q02_1}{Enthusiast Pre-Score - Test 03 (23 Nov 2020)-img/p2_math_s3_q02_1.png}{156.5}{0}
\raggedcolumns
\makeatletter\@beginparpenalty=10000\makeatother
\newcolumntype{L}[1]{>{\raggedright\arraybackslash}p{#1}}
\begin{document}
\begin{center}{\Large\bfseries Enthusiast Pre-Score --- Test 03 (23 Nov 2020)}\end{center}
\begin{multicols}{2}
\raggedcolumns
\papersec{1}
\partsec{1}{PHYSICS}
\sect{SECTION-I(i)}
\qn{1} Two long conducting rectangular strips are placed parallel to each other at a distance d as shown. Strip A is connected to a non-conducting spring of spring constant k. Strip B is fixed. Given $\ell\gg b\gg d$. A current i flows uniformly along the length $\ell$ of each strip and the currents in the strips are anti-parallel to each other. The initial deformation in the spring keeps the system in static equilibrium. Now at t = 0, the currents in both strips are halved and made parallel to each other suddenly. (Neglect gravity and assume force on A is due to B and spring only and consider A and B as large sheets and neglect end effects)\par
\qfig{p1_phy_s1_q01_1}
\begin{qoptions}
\opt{A}{the instantaneous force on strip A at t = 0$^+$ is given by $\dfrac{5\mu_0i^2\ell}{8b}$.}
\opt{B}{the instantaneous force on strip A at t = 0$^+$ is given by $\dfrac{3\mu_0i^2\ell}{8b}$.}
\opt{C}{Elongation in spring at the instant t = 0$^+$ is given by $\dfrac{\mu_0i^2\ell}{2bk}$.}
\opt{D}{Elongation in spring at the instant t = 0$^+$ is given by $\dfrac{\mu_0i^2\ell}{4bk}$}
\end{qoptions}
\qn{2} The figure-2, figure-3, figure-4 and figure-5 represents the graphs of the source voltage V, resistor voltage V$_R$, inductor voltage V$_L$ and capacitor voltage V$_C$ as a function of time t of the series circuit with impedance 500$\Omega$ as shown in the figure-1 $\left(t_0=\dfrac{\pi}{20}\text{ms}\right)$. Choose correct option(s).\par
\qfig{p1_phy_s1_q02_2}
\qfig{p1_phy_s1_q02_1}
\qfig{p1_phy_s1_q02_3}
\qfig{p1_phy_s1_q02_4}
\qfig{p1_phy_s1_q02_5}
\begin{qoptions}
\opt{A}{R = 400 $\Omega$, C = 0.5 $\mu$F and L = 50 mH}
\opt{B}{R = 300 $\Omega$, C = 0.5 $\mu$F and L = 60 mH}
\opt{C}{V$_{10}$ = 30V}
\opt{D}{I = 0.1A cos (10000 t) Amp and current lag the source voltage}
\end{qoptions}
\qn{3} A straight current carrying wire extends along z-axis from A (0, 0, $-$2) to B (0,0,2), A current of 6A flows from A to B. Magnetic field intensity at the point $(\sqrt2\text{m},\sqrt2\text{m},0\text{m})$ will be :-
\begin{qoptions}
\opt{A}{$3\sqrt2\times10^{-7}$T}
\opt{B}{3 $\times$ 10$^{-7}$ T}
\opt{C}{Directed along $(\hat i-\hat j)$}
\opt{D}{Directed along $(-\hat i+\hat j)$}
\end{qoptions}
\qn{4} A standing wave is generated on a string of length 10m such that string vibrates in 10$^{\text{th}}$ harmonic. Frequency of the source used was 1000 Hz. It was found that at t = 0, all the particles were at half of their amplitude and particle at x = 0.5m was located in +ve y. If amplitude of antinode is A \& string's one end is fixed, choose the possible wave equation treating x = 0 to be fixed end.
\begin{qoptions}
\opt{A}{$y=A\sin(\pi x+\pi)\sin\left(2000\pi t+\dfrac{7\pi}6\right)$}
\opt{B}{$y=A\sin\pi x\,\sin\left(2000\pi t+\dfrac\pi6\right)$}
\opt{C}{$y=A\sin\pi x\,\sin\left(2000\pi t+\dfrac{5\pi}6\right)$}
\opt{D}{$y=A\sin(\pi x+\pi)\sin\left(2000\pi t+\dfrac\pi6\right)$}
\end{qoptions}
\qn{5} The figure shows a section of a tube held fixed by an external force, and liquid of density $\rho$ flowing through it. There is uniform flow across any cross-section of the flow. The inlet area is A$_1$ and speed is v$_1$ and P$_1$ is the pressure of liquid at inlet and outlet area is A$_2$ and speed is v$_2$ and pressure P$_2$. The flow is incompressible. Which of the following option is/are correct ? Apart from section MNOP of the tube, other parts are perfectly cylindrical with horizontal axis. (Ignore atmospheric pressure).\par
\qfig{p1_phy_s1_q05_1}
\begin{qoptions}
\opt{A}{Energy conservation is valid in an inertial frame going right with speed v$_2$}
\opt{B}{Bernoull's equation is valid in an inertial frame going right with speed v$_2$}
\opt{C}{If the external force is removed, the tube accelerates towards left.}
\opt{D}{If the external force is removed, the tube accelerates towards right.}
\end{qoptions}
\qn{6} Capacitor with a capacitance C and an inductor L are connected to each other, and in the resulting circuit oscillations occur. At the moment when the capacitor voltage was U, and current I flowed through the coil, we connected a resistor of resistance R in parallel to the coil. How much heat H is dissipated in the resistor ? What modulus of charge Q flows through coil starting from now on ?
\begin{qoptions}
\opt{A}{$H=\dfrac{CU^2}{2}+\dfrac{LI^2}{2}$}
\opt{B}{$Q=\dfrac{LI}{R}$}
\opt{C}{$Q=\dfrac{LI}{R}-CU$}
\opt{D}{$Q=\dfrac{LI}{R}+CU$}
\end{qoptions}
\sect{SECTION-I(ii)}
\paraq{7}{8}\textbf{REFLECTION}. A wave pulse on a string has the dimensions shown in figure at t = 0. The wave speed is 40 cm/s.\par
\qfig{p1_phy_s2_para7_8_1}
\qn{7} If point O is a fixed end, draw the shape of the string at t = 35 ms.
\begin{qoptions}
\opt{A}{\qfigi{p1_phy_s2_q07_A}}
\opt{B}{\qfigi{p1_phy_s2_q07_B}}
\opt{C}{\qfigi{p1_phy_s2_q07_C}}
\opt{D}{\qfigi{p1_phy_s2_q07_D}}
\end{qoptions}
\qn{8} Repeat for the case when point A is the free end and t = 30 ms
\begin{qoptions}
\opt{A}{\qfigi{p1_phy_s2_q08_A}}
\opt{B}{\qfigi{p1_phy_s2_q08_B}}
\opt{C}{\qfigi{p1_phy_s2_q08_C}}
\opt{D}{\qfigi{p1_phy_s2_q08_D}}
\end{qoptions}
\paraq{9}{10}A cylindrical capacitor with external radius R, internal radius R$-$d (d $\ll$ R), length $l$ and mass M hangs on an insulating cord in a region where there is a homogenous, vertical magnetic field of strength B. It can rotate freely as a whole around its vertical axis, but is constrained, so that it can not move horizontally. The capacitor is charged and there is a voltage difference V between its plates. (Neglect gravity)\par\par
\qfig{p1_phy_s2_para9_10_1}
\qn{9} Charge on capacitor is given by
\begin{qoptions}
\opt{A}{$\epsilon_0\dfrac{2\pi R/V}{d}$}
\opt{B}{$\dfrac{\epsilon_0\pi R/V}{2d}$}
\opt{C}{$\dfrac{\epsilon_0\pi R/V}{3d}$}
\opt{D}{$\dfrac{\epsilon_0\pi R/V}{4d}$}
\end{qoptions}
\qn{10} If without being mechanically disturbed, the capacitor is discharged through an internal radial wire then find maximum angular velocity of capacitor
\begin{qoptions}
\opt{A}{$\dfrac{\pi\epsilon_0VBl}{M}$}
\opt{B}{$\dfrac{2\pi\epsilon_0VBl}{M}$}
\opt{C}{$\dfrac{3\pi\epsilon_0VBL}{M}$}
\opt{D}{$\dfrac{4\pi\epsilon_0VBL}{M}$}
\end{qoptions}
\sect{SECTION-II}
\qn{1} The coil in a moving coil Galvanometer experiences torque proportional to current passed through it. If a steady current i is passed through it the deflection of the coil is found to be 90$^\circ$. Now the steady current is switched off and a charge Q is suddenly passed through the coil. If the coil has N turns of area A each and its moment of inertia is I about the axis it is going to rotate and placed in uniform magnetic field B then the maximum angle through which it deflects upon passing Q is $Q\sqrt{\dfrac{\pi NAB}{k\,i\,I}}$. Find the value of k.(Assume damping is small and it takes large number of oscillations for coil to eventually stop)
\qn{2} In an L--R series circuit a sinusoidal voltage v = v$_0$ sin $\omega$t is applied. It is given that L = 35 mH, R = 11 $\Omega$, V$_{rms}$ = 220 V, $\dfrac{\omega}{2\pi}=50$Hz and $\pi=\dfrac{22}{7}$. The amplitude of current is 4n A. Find n.
\qn{3} A charge particle is projected in gravity free space from origin with $\vec v=(2\hat i+2\hat j+2\sqrt3\,\hat k)$ m/s. There is uniform magnetic field $\vec B=(2\hat j+2\sqrt3\,\hat k)$ T and uniform electric field $\vec E=(4\sqrt3\,\hat j-4\hat k)\dfrac NC$ in space. What is the distance travelled by charge in time $\sqrt5$ sec from start (in meter) ? Take charge = 1mC \& mass = 20 g
\qn{4} Figure represent a infinite long wire with current I. There are two current loop of same small width $\ell$ subtending equal angle at wire ($\ell\ll r$). Neglect interaction of two current loop. What is ratio of magnetic torque due to wire on bigger loop to torque on smaller loop ?\par
\qfig{p1_phy_s3_q04_1}
\qn{5} Suddenly uniform magnetic field B$_0$ ($-\hat k$) is applied in the circular gravity free region having radius `4a'. Initially field was zero and charge (q, m) was at rest at (a, 0). Now once field is setup find the radius of circular path (in meter)\par
Given : a = 4.00 m.\par
\qfig{p1_phy_s3_q05_1}
\qn{6} A uniform meter scale is balanced on a fixed semicircular cylinder of radius 30 cm as shown. One end of scale is slightly depressed and released. Time period of oscillation (in sec) is (There is no slipping between scale and surface of cylinder)\par
\qfig{p1_phy_s3_q06_1}
\qn{7} A rod of mass `m' and length $\ell$ is connected by two springs of spring constant k$_1$ and k$_2$, so that is horizontal at equilibrium. What is the natural frequency of the system ?\par
(Given m = 3 kg, b = 1 m, $\ell$ = 2m, k$_1$ = 10 N/m, k$_2$ = 20 N/m and take $\pi=\sqrt{10}$)\par
\qfig{p1_phy_s3_q07_1}
\qn{8} Two spheres A and B of the same mass m and the same radius are placed on a rough hoizontal surface A is a uniform hollow sphere and B is uniform solid sphere A and B can roll without sliding on the floor. They also tied centrally to a light spring of spring constant k as shown in figure. A and B are released when the extension in the spring is x$_0$ the amplitude of SHM of the spheres A is $\dfrac{x.x_0}{46}$ then x is :-\par
\qfig{p1_phy_s3_q08_1}
\partsec{2}{CHEMISTRY}
\sect{SECTION-I(i)}
\qn{1} Consider different complex where M$^{+2}$ represent Ni$^{+2}$ (X, Y are neutral monodentated ligand) and Find \textbf{CORRECT} statement(s).
\begin{qoptions}
\opt{A}{If [MX$_2$Y$_2$]$^{+2}$ is paramagnetic then it has planar geometry at metal}
\opt{B}{If [MX$_2$Y$_2$]$^{+2}$ is diamagnetic then it has two stereoisomer}
\opt{C}{[MX$_6$]$^{+2}$ is paramagnetic with sp$^3$d$^2$ hybridisation}
\opt{D}{Number of valence electrons on metal [MX$_4$]$^{+2}$ and [MY$_6$]$^{+2}$ are same}
\end{qoptions}
\qfullbegin
\qn{2} \qfig{p1_chem_s1_q02_1}
\qfullsplit
Select \textbf{CORRECT} statements.
\begin{qoptions}
\opt{A}{[A]$_{aq.}$ can be solution of CuO in dil. H$_2$SO$_4$}
\opt{B}{Addition of [C] to [A] is not a redox change}
\opt{C}{Complex in colourless solution is diamagnetic \& tetrahedral}
\opt{D}{Solution (C) can be NH$_4$OH}
\end{qoptions}
\qfullend
\qn{3} For a zero order reaction: A $\to$ 2B, [A] = 0.4M and [B] = 0 at t = 0. If at t = 25 min, [B] = 0.5M, then the correct information(s) about the reaction, with respect to initial condition, is/are:
\begin{qoptions}
\opt{A}{t$_{100\%}$ = 40 min}
\opt{B}{T$_{av}$ = 20 min}
\opt{C}{t$_{50\%}$ = 20 min}
\opt{D}{t$_{10\%}$ = 4 min}
\end{qoptions}
\qn{4} The correct information(s) about nuclear processes is/are:
\begin{qoptions}
\opt{A}{In alpha decay, two protons and two neutrons come out from the nucleus, in succession.}
\opt{B}{An independent neutron may convert spontaneously into proton.}
\opt{C}{An independent proton may convert spontaneously into neutron.}
\opt{D}{In a positron decay, a proton spontaneously converts into neutron.}
\end{qoptions}
\qn{5} In which of the following reaction(s), 2$^\circ$ alcohol is not obtained as major product.\par
\qfig{p1_chem_s1_q05_A}
\qfig{p1_chem_s1_q05_B}
\qfig{p1_chem_s1_q05_C}
\qfig{p1_chem_s1_q05_D}
\qn{6} \qfig{p1_chem_s1_q06_1}
But-1-ene is obtained as major product when :
\begin{qoptions}
\opt{A}{(G) is $-$F}
\opt{B}{Base is $\mathrm{CH_3\text{--}C(CH_3)_2\text{--}ONa}$}
\opt{C}{(G) is $-\overset{\oplus}{\mathrm N}\mathrm{Me_3}$}
\opt{D}{(G) is $-$Cl \& base is alc.KOH, $\Delta$}
\end{qoptions}
\sect{SECTION-I(ii)}
\paraq{7}{8}In given reaction sequence\par
\qfig{p1_chem_s2_para7_8_1}
\qn{7} Product `R' is :
\begin{qoptions}
\opt{A}{\qfigi{p1_chem_s2_q07_A}}
\opt{B}{\qfigi{p1_chem_s2_q07_B}}
\opt{C}{\qfigi{p1_chem_s2_q07_C}}
\opt{D}{\qfigi{p1_chem_s2_q07_D}}
\end{qoptions}
\qn{8} Product R can't obtained as a major product in reaction :\par
\qfig{p1_chem_s2_q08_A}
\qfig{p1_chem_s2_q08_B}
\qfig{p1_chem_s2_q08_C}
\qfig{p1_chem_s2_q08_D}
\qn{9} The complex first denoted by $\beta$-[PtCl$_2$(NH$_3$)$_2$] was identified as the trans-isomer. (The cis isomer was denoted by $\alpha$) It reacts slowly with solid Ag$_2$O to produce [Pt(NH$_3$)$_2$(OH$_2$)$_2$]$^{2+}$. This complex does not react with C$_2$O$_4^{2-}$ to give a chelated complex.\par
The `third isomer' (neither $\alpha$ nor $\beta$) of emperical composition PtCl$_2\cdot$2NH$_3$ is an insoluble solid which, when heated with aq. AgNO$_3$, gives a solution containing [Pt(NH$_3$)$_4$](NO$_3$)$_2$ and a new solid phase of composition Ag$_2$[PtCl$_4$].\par
Which of the following structure of complex is \textbf{NOT} observed in any platinum complex compound in paragraph.\par
($\bigcirc$ = Pt$^{+2}$)
\begin{qoptions}
\opt{A}{\qfigi{p1_chem_s2_q09_A}}
\opt{B}{\qfigi{p1_chem_s2_q09_B}}
\opt{C}{\qfigi{p1_chem_s2_q09_C}}
\opt{D}{\qfigi{p1_chem_s2_q09_D}}
\end{qoptions}
\qn{10} What is the magnetic nature, hybridization state of central metal ion, hybridisation of nitrogen in co-ordination sphere \& hybridization of nitrogen in ionisation sphere respectively in \textbf{[Pt(NH$_3$)$_4$](NO$_3$)$_2$}.
\begin{qoptions}
\opt{A}{(diamagnetic) dsp$^2$, sp$^3$, sp$^3$}
\opt{B}{(diamagnetic) dsp$^2$, sp$^3$, sp$^2$}
\opt{C}{(diamagnetic) sp$^3$, sp$^3$, sp$^2$}
\opt{D}{(paramagnetic) dsp$^2$, sp$^3$, sp$^2$}
\end{qoptions}
\sect{SECTION-II}
\qn{1} An element \textbf{`E'} of 4$^{\text{th}}$ period of periodic table which exist in B.C.C. lattice structure, have high melting point and large heat of atomisation. Element \textbf{`E'} have \textbf{E$^{+2}$/E} electrode potential = $-$1.18 V and \textbf{E$^{+3}$/E$^{+2}$} electrode potential = $-$0.26 V.\par
Element \textbf{`E'} produces violet coloured aquated ion [E(H$_2$O)$_6$]$^{+2}$, which can reduce H$_2$O with the liberation of H$_2$, thus the violet solution soon turns green through oxidation to [E(H$_2$O)$_6$]$^{+3}$ paramagnetic ion.\par
If atomic number of \textbf{E} = `x' and highest oxidation state of \textbf{E} = +`y' then \textbf{find (x + y)}.
\qn{2} Hexadentate ligand forms very stable complexes (usually octahedral structures) with most of the transition metals. Find number of five member chelate ring(s) which must involve O--Co--N bond angle in $\overset{\text{III}}{\mathrm{Co}}\mathrm{(EDTA)}^-$.\par
(where EDTA$^{4-}$ is ethylendiamine tetraacetate ion)
\qfullbegin
\qn{3} Total number of reactions correctly match with their major products are ?\par
\qfig{p1_chem_s3_q03_1}
\qfullend
\qn{4} \qfig[0.8]{p1_chem_s3_q04_1}
Number of equivalents of Grignard reagent (x) used in the reaction is
\qn{5} Liquids `A' and `B' form an ideal solution. At 50$^\circ$C, $\mathrm{P_A^\circ}$ = 0.25 bar and $\mathrm{P_B^0}$ = 0.50 bar. All the vapours at equilibrium above a solution containing liquids `A' and `B' in 4 : 1 mole ratio at 50$^\circ$C, is collected in an another vessel and condensed. The condensate is heated to 50$^\circ$C and all the vapours at equilibrium is again collected in a third vessel. Similar process is repeated several times. What is the mole-fraction of `A' in the 4th condensate ?
\qn{6} A sample contains two ideal gases, `A' and `B' in $\sqrt2:\sqrt3$ molar ratio. If the ratio of molar masses, M$_A$ : M$_B$ = 1 : 2 and the ratio of their radii, r$_A$ : r$_B$ = 2 : 3, the ratio of collision frequencies, Z$_{AA}$ : Z$_{AB}$ is x : 1.00. The value of \textbf{`x'} is
\qn{7} 40 gm of a sample of caustic soda containing NaOH, Na$_2$CO$_3$ and inert impurity is dissolved in water to prepare 1.0 litre solution. A 25 ml portion of this solution required 24 ml 1.0 M-HCl for complete neutralisation. To another 25 ml solution, excess of BaCl$_2$ is added, and after filtration, the resulting solution required 22 ml HCl of same strength to reach end point. Calculate the mass percent of Na$_2$CO$_3$ in the original sample.
\qn{8} Approximately how many oxygen molecules (in the order 10$^6$) arrive each second at each mitochondrion of an active person with a mass of 84 kg? The following data are available: Oxygen consumption is about 48 ml per minute per kilogram of body weight, measured at 300K and 1 atm. In an adult, there is about 1.6 $\times$ 10$^{10}$ cells per kg body mass and each cell contains about 800 mitochondria. (R = 0.08 L-atm/K-mol, NA = 6 $\times$ 10$^{23}$).
\partsec{3}{MATHEMATICS}
\sect{SECTION-I(i)}
\qn{1} Consider a pair of straight lines P : $x^2+y^2-2xy-5x+5y+6=0$
\begin{qoptions}
\opt{A}{Area bounded by P and co-ordinate axes is $\dfrac32$ sq. units}
\opt{B}{Area bounded by P and co-ordinate axes is $\dfrac52$ sq. units}
\opt{C}{The four intersection points of P with co-ordinate axes are concyclic}
\opt{D}{The four intersection points of P with co-ordinate axes are non-concyclic}
\end{qoptions}
\qn{2} Figure shows a circle $x^2+y^2=a^2$ whose centre `O' is at origin. Line y = mx + c drawn parallel to AB cuts circle at M and N \& co-ordinate axes at P \& Q as shown. If MP = $\sqrt{56}$ \& PN = 12, then -\par
\qfig{p1_math_s1_q02_1}
\begin{qoptions}
\opt{A}{$c^2=100+12\sqrt{56}$}
\opt{B}{$c^2=100-12\sqrt{56}$}
\opt{C}{m = $-$1}
\opt{D}{$a^2$ = 100}
\end{qoptions}
\qn{3} Five digit numbers are formed using digits 1,2,3,4,5,6. The number of such numbers are -
\begin{qoptions}
\opt{A}{7500, if only two non-adjacent digits are same}
\opt{B}{2160, if only two non-adjacent digits are same}
\opt{C}{1440, if only two adjacent digits are same}
\opt{D}{3600, if only two digits are same}
\end{qoptions}
\qn{4} Coordinates of point A is (1,2) and P($\alpha$,$\beta$) lies on the line L$_1$ : 2x + 3y + 5 = 0. Point B($\gamma$,$\delta$) is point of intersection of L$_2$ : 3x $-$ 2y = 0 \& L$_3$ : $\lambda$x + 2y = 10, $\lambda\in$ R$^+$ such that $\gamma$ \& $\delta$ are positive integers.\par
Let s = $\alpha+\beta+\gamma+\delta+\lambda$ , then which of following is/are correct-
\begin{qoptions}
\opt{A}{If $|$PA $-$ PB$|$ is minimum, then s = 11}
\opt{B}{If $|$PA $-$ PB$|$ is maximum, then 5s = 24}
\opt{C}{If $|$PA + PB$|$ is minimum, then 31s = 156}
\opt{D}{If $|$PA + PB$|$ is minimum, then 5s = 24}
\end{qoptions}
\qn{5} If H $\equiv$ (3, 4) and O $\equiv$ (1, 2) are orthocentre and circumcentre of $\Delta$PQR respectively and equation of side PQ is x $-$ y + 7 = 0, then -
\begin{qoptions}
\opt{A}{equation of circumcircle of $\Delta$PQR is $(x-1)^2+(y-2)^2=80$}
\opt{B}{equation of circumcircle of $\Delta$PQR is $(x-1)^2+(y-2)^2=100$}
\opt{C}{centroid of $\Delta$PQR is $\left(\dfrac53,\dfrac83\right)$}
\opt{D}{midpoint of PQ is ($-$2, 5)}
\end{qoptions}
\qn{6} If $T_r=\left|\begin{array}{ccc}r-1&r-2&r-6\\2r-4&2r-6&2r-11\\3r-9&3r-12&3r-18\end{array}\right|$, then -
\begin{qoptions}
\opt{A}{$\displaystyle\sum_{r=0}^{5}T_r=36$}
\opt{B}{$\displaystyle\sum_{r=1}^{5}T_r=30$}
\opt{C}{$\displaystyle\prod_{r=0}^{5}T_r=6^6$}
\opt{D}{$\displaystyle\prod_{r=1}^{5}T_r=6^5$}
\end{qoptions}
\sect{SECTION-I(ii)}
\qn{7} $C_1:x^2+y^2=25$, $C_2:x^2+y^2=16$ are the given circles A($-$3,4) is a fixed point on C$_1$ and B is variable point on C$_2$. Let AB meet C$_2$ again at C so that B lies on the line segment AC. Then\par
Maximum value of $(BC)^2$ is -
\begin{qoptions}
\opt{A}{100}
\opt{B}{64}
\opt{C}{36}
\opt{D}{9}
\end{qoptions}
\qn{8} If $(BC)^2$ is maximum when the coordinates of C are-
\begin{qoptions}
\opt{A}{$\left(\dfrac{12}5,-\dfrac{16}5\right)$}
\opt{B}{$\left(\dfrac{12}5,\dfrac{16}5\right)$}
\opt{C}{$\left(-\dfrac{12}5,\dfrac{16}5\right)$}
\opt{D}{$\left(-\dfrac{12}5,-\dfrac{16}5\right)$}
\end{qoptions}
\qn{9} Consider a word EARTHQUAKE.\par
The number of words in which EA as well as AE appear but they are not consecutive-
\begin{qoptions}
\opt{A}{$\dfrac{10!}{2!2!}$}
\opt{B}{7(7!)}
\opt{C}{6(7!)}
\opt{D}{8!}
\end{qoptions}
\qn{10} Number of selection of 5 letters from the given word in which two are alike and three are different is-\par
(for example AARTE)
\begin{qoptions}
\opt{A}{40}
\opt{B}{30}
\opt{C}{35}
\opt{D}{70}
\end{qoptions}
\sect{SECTION-II}
\qn{1} A variable line ax + by + c = 0, where a, b, c are in A.P., is normal to a circle $(x-\alpha)^2+(y-\beta)^2=\gamma$, which is orthogonal to circle x$^2$ + y$^2$$-$ 4x$-$ 4y$-$1 = 0. The value of $\left(\dfrac{\alpha+\beta+\gamma}{4}\right)$ is equal to
\qn{2} The equation of straight line passing through the point of intersection of family of line x(3 + 4$\lambda$) + y(4 + 3$\lambda$) + (1 + 6$\lambda$) = 0 which lies at maximum distance from the point (1, 1) is x + by + c = 0, then value of (b + c) is
\qn{3} If sum of areas of all the concave quadrilaterals formed by points A(0,0), $B\left(-1,\dfrac{16}3\right)$, C(3,4) and D(1,3) be $\lambda$, then $\left(\dfrac\lambda8\right)$ equals
\qn{4} Let $\lambda$ be the positive factor of 1200. If the number of positive integral solution of x$_1$x$_2$x$_3$ = $\lambda$ (for all possible values of $\lambda$) is P, then number of divisors of P equals
\qn{5} Two circles of radii r$_1$ and r$_2$ touch each-other externally such that $r_2=\left(3+2\sqrt2\right)r_1$. If 5x$^2$ + 3xy + $\lambda$y$^2$ = 0 is a pair of direct common tangent between them, then the value of $\dfrac{|\lambda|}{4}$ is
\qn{6} The value of `$\lambda$' if $ax^4+bx^3+cx^2+50x+d=\left|\begin{array}{ccc}x^3-14x^2&-x&3x+\lambda\\4x+1&3x&x-4\\-3&4&0\end{array}\right|$ is -
\qn{7} Let a$_n$ denotes the number of sequences of length n formed by using the digits 0,1 such that there are no two consecutive 1 in the sequence, then the value of $\dfrac15\left(\left(\dfrac{a_5a_6-a_5^2}{a_4}\right)-5\left(\dfrac{a_n-a_{n-1}}{a_{n-2}}\right)\right)$ is equal to\par
(where 010 , 001,100, 101, all are sequences of length three)
\qn{8} If system of equations\par
$cx+y+1=0$\par
$x+cy+2=0$\par
$x+y+1=0$\par
is conistent and the sum of the values of c is equal to $\lambda$ then $\dfrac\lambda5$ equals
\papersec{2}
\partsec{1}{PHYSICS}
\sect{SECTION-I(i)}
\qn{1} A hemispherical container of mass M is placed upside down on a smooth horizontal surface. Through a small hole at the top of container, water is poured into it as shown in the figure. Exactly when the container gets fill (completely filled), water starts leaking from between the table and the edge of the container. The density of water is $\rho$ and radius of hemi-spherical container is R. (Neglect the atmospheric pressure and effect of surface tension) choose the correct statement(s).\par
\qfig{p2_phy_s1_q01_1}
\begin{qoptions}
\opt{A}{The force exerted on the container due to liquid, when half filled (upto height R/2) is $\dfrac1{24}\pi R^3\rho g$}
\opt{B}{The force exerted on the container due to liquid, when half filled (upto height R/2) is $\dfrac78\pi R^3\rho g$}
\opt{C}{The mass `M' of container is $\dfrac{2\pi R^3\rho}{3}$}
\opt{D}{The mass `M' of container is $\dfrac{\pi R^3\rho}{3}$}
\end{qoptions}
\qn{2} The circuit shown consists of two unknown resistances $R_1$ and $R_2$, a variable reactance $X_L$ and an AC voltage source. The peak values of currents supplied by the source are 9 A and 3 A when the variable reactance is made 0 and infinity respectively. Now the reactance is so adjusted that the overall power factor of the circuit has its minimum value. Voltage of the source indicated in the figure is the peak voltage. Which of the following conclusion can you make?\par
\qfig{p2_phy_s1_q02_1}
\begin{qoptions}
\opt{A}{Peak current supplied by the source is $3\sqrt3$ A 30$^\circ$ lagging behind voltage..}
\opt{B}{The required value of the reactance is $4\sqrt3\ \Omega$.}
\opt{C}{Total power dissipated in the circuit is 54 W}
\opt{D}{The power factor of circuit is $\dfrac{\sqrt3}2$}
\end{qoptions}
\qn{3} Time dependent magnetic field is set up in a circular region as shown in figure. $\dfrac{dB}{dt}$ is positive. Choose \textbf{CORRECT} option(s).\par
\qfig{p2_phy_s1_q03_1}
\begin{qoptions}
\opt{A}{$\displaystyle\int_Q^R\vec E\cdot\vec{dl}$ is positive; for straight line path QR.}
\opt{B}{$\displaystyle\int_P^Q\vec E\cdot\vec{dl}$ is negative; for straight line path PQ.}
\opt{C}{$\displaystyle\int_R^P\vec E\cdot\vec{dl}$ is negative; for curve arc path RQP.}
\opt{D}{$\displaystyle\int_P^Q\vec E\cdot\vec{dl}=0$ ; for straight line path PQ.}
\end{qoptions}
\qn{4} A square rigid loop of dimension `a' meter is placed in a region of uniform magnetic field B, as shown. Half of the loop is inside \& other half is outside the field. Self inductance of the loop is L = 100 henry \& the resistance of the loop is R = 50$\Omega$. A capacitor is also connected in loop in one of the side of loop having capacitance C = 1 $\mu$F. At t = 0 sec, external force (agent) starts moving loop according to equation $x=\dfrac a2\sin(\omega_0t)$, where $\omega_0=100$ rad/s.\par
\qfig{p2_phy_s1_q04_1}
\begin{qoptions}
\opt{A}{Current starts with anticlockwise direction, after some time its direction may change.}
\opt{B}{Peak value of current in loop is $\left(\dfrac{B\omega_0a^2}{2R}\right)$.}
\opt{C}{Average power delivered by external agent in one cycle is $\left(\dfrac{B^2\omega_0^2a^2}{8R}\right)$ once steady state is reached.}
\opt{D}{If angular frequency of oscillation is increased from $\omega_0$ then average power delivered by external agent will decrease.}
\end{qoptions}
\qn{5} Figure shows a spherical body floating on liquid surface. Line AB represents diameter of sphere. Assume liquid and material of body makes a non wetting combination with contact angle = 143$^\circ$ at the circumference of sphere passing through AB. surface tension = 0.05 N; density of liquid = 4 $\times$ 10$^3$ kg/m$^3$; material of radius of sphere = 2 mm, g = 10m/s$^2$. The extrapolated line AB is at a depth h below the horizontal part of free surface.\par
\qfig{p2_phy_s1_q05_1}
\begin{qoptions}
\opt{A}{h = 1 mm}
\opt{B}{h = 1.2 mm}
\opt{C}{The density of material of sphere is 5 $\times$ 10$^3$ kgm$^{-3}$}
\opt{D}{The density of material of sphere is 4.5 $\times$ 10$^3$ kgm$^{-3}$}
\end{qoptions}
\qn{6} To the ``bridge'' of capacitors (Fig.), a battery with a voltage of U was connected. Then it was removed, and between points A and B we connected a coil of inductance L. Find the maximum current I through the coil. Find also the total charge Q flowing through the coil, and the amount of heat H dissipated in it a very-very long time. The resistance of the connecting wires is zero, the resistance R of the wire with which the coil is wound is considered very small but not zero.\par
\qfig{p2_phy_s1_q06_1}
\begin{qoptions}
\opt{A}{$I=\sqrt{\dfrac{3C}{8L}}U$}
\opt{B}{$Q=\dfrac{3CU}{4}$}
\opt{C}{$H=\dfrac{3CU^2}{16}$}
\opt{D}{When the current in coil is maximum for the first time, the charge on each of the capacitors is zero.}
\end{qoptions}
\sect{SECTION-I(ii)}
\qn{7} \textbf{List--I (System)} \quad \textbf{List--II (Maximum amplitude of block so that string attached to block does not slack)}\par
(P)\par
\qfig{p2_phy_s2_q07_1}
(Q)\par
\qfig{p2_phy_s2_q07_2}
(R)\par
\qfig{p2_phy_s2_q07_3}
(S)\par
\qfig{p2_phy_s2_q07_4}
(1) $\dfrac{mg}{2k}$ \quad (2) $\dfrac{4mg}{k}$ \quad (3) $\dfrac{mg}{8k}$ \quad (4) $\dfrac{mg}{k}$
\begin{qoptions}
\opt{A}{P $\to$ 4; Q $\to$ 3; R $\to$ 1; S $\to$ 2}
\opt{B}{P $\to$ 4; Q $\to$ 1; R $\to$ 3; S $\to$ 2}
\opt{C}{P $\to$ 4; Q $\to$ 3; R $\to$ 2; S $\to$ 1}
\opt{D}{P $\to$ 4; Q $\to$ 2; R $\to$ 3; S $\to$ 1}
\end{qoptions}
\qfullbegin
\qn{8} List-I shows charges projected in given field, list-II shows path of respective charges.\par
\qfig{p2_phy_s2_q08_1}
\qfullsplit
\begin{qoptions}
\opt{A}{P $\to$ 2; Q $\to$ 1; R $\to$ 4; S $\to$ 3}
\opt{B}{P $\to$ 3; Q $\to$ 1; R $\to$ 4; S $\to$ 2}
\opt{C}{P $\to$ 4; Q $\to$ 2; R $\to$ 1; S $\to$ 3}
\opt{D}{P $\to$ 3; Q $\to$ 1; R $\to$ 2; S $\to$ 4}
\end{qoptions}
\qfullend
\qn{9} In the figure below two strings of different linear mass density $\mu_1$ and $\mu_2$ and length $\ell_1$ and $\ell_2$ are connected to each other. The string passes over the ideal pulley and connected to weight `W'. When excited with a tuning fork the string vibrates such that the joint between the two wires does not move at all. If N$_1$ and N$_2$ are the number of harmonics in the first and second wire respectively, match the entries in List-I and List-II.\par
\qfig{p2_phy_s2_q09_1}
\noindent{\footnotesize\setlength{\tabcolsep}{2pt}\begin{tabular}{|L{.8cm}|L{2.5cm}|L{.8cm}|L{3.3cm}|}\hline
\multicolumn{2}{|c|}{List-I} & \multicolumn{2}{c|}{List-II}\\\hline
(P) & $\dfrac{\mu_2}{\mu_1}=4:1$ & (1) & $\ell_1=3\ell_2$, N$_1$ : N$_2$ = 1 : 1\\\hline
(Q) & $\dfrac{\mu_1}{\mu_2}=1:9$ & (2) & $\ell_1:\ell_2=1:6$, N$_1$ : N$_2$ = 1 : 3\\\hline
(R) & $\dfrac{\mu_2}{\mu_1}=1:4$ & (3) & $\ell_1:\ell_2=3:1$, N$_1$ : N$_2$ = 3 : 2\\\hline
(S) & $\dfrac{\mu_1}{\mu_2}=4:9$ & (4) & $2\ell_1=\ell_2$, N$_1$ : N$_2$ = 1 : 3\\\hline
\end{tabular}}
\begin{qoptions}
\opt{A}{P $\to$ 4; Q $\to$ 1; R $\to$ 3; S $\to$ 2}
\opt{B}{P $\to$ 2; Q $\to$ 1; R $\to$ 4; S $\to$ 3}
\opt{C}{P $\to$ 3; Q $\to$ 4; R $\to$ 2; S $\to$ 1}
\opt{D}{P $\to$ 3; Q $\to$ 1; R $\to$ 2; S $\to$ 4}
\end{qoptions}
\qn{10} There is a very long vertical cylindrical tube containing a viscous liquid. A ball is released in the liquid with initial conditions given in column - I. The density of ball) may be greater or less than that of liquid. From its initial position, the ball can move upward or downward for a very long time. Consider upward direction as positive for v-t graphs in column - II.\par
\textbf{List-I}\par
(P) The ball is given initial velocity in upward direction in uniform gravity downward\par
(Q) The ball is released with zero initial velocity in uniform gravity downward\par
(R) The ball is given initial velocity u in upward direction and the tube is placed in gravity free space\par
(S) The ball is released with zero initial velocity but the liquid medium is moving upward with velocity uo and tube is placed in gravity free space\par
\textbf{List-II}\par
(1)\par
\qfig{p2_phy_s2_q10_1}
(2)\par
\qfig{p2_phy_s2_q10_2}
(3)\par
\qfig{p2_phy_s2_q10_3}
(4)\par
\qfig{p2_phy_s2_q10_4}
\begin{qoptions}
\opt{A}{P $\to$ 2; Q $\to$ 3; R $\to$ 4; S $\to$ 1}
\opt{B}{P $\to$ 3; Q $\to$ 4; R $\to$ 1; S $\to$ 2}
\opt{C}{P $\to$ 1,3; Q $\to$ 2; R $\to$ 4; S $\to$ 2}
\opt{D}{P $\to$ 2,3; Q $\to$ 4; R $\to$ 3; S $\to$ 1}
\end{qoptions}
\sect{SECTION-II}
\qn{1} An ideal capacitor and an near ideal (very small resistance) copper wire loop inductor are connected in series with an AC source, of rms voltage 200 volt \& variable frequency. With a frequency (f$_1$ \& f$_2$) of AC source power drawn by circuit is half of maximum power drawn. At frequency $\dfrac{f_1+f_2}{2}$, peak current in circuit is $50\sqrt2$ A. Now this inductor is connected with DC cell of EMF 5V and internal resistance 1$\Omega$. Find steady state current in inductor in ampere.
\qn{2} A square loop carrying a current I lies in the y-z plane with its vertices at points (0, a, a), (0, $-$a, a), (0, a, 3a) and (0, $-$a, 3a). Find the magnitude of net magnetic flux due to the current loop through the first quadrant in the x-y plane.\par
Take a = 2m, I = 3A, ln 2 = 0.7, ln3 = 1.1. If the answer is $\phi$ Wb, fill $\phi\times10^7$ in OMR sheet.\par
\qfig{p2_phy_s3_q02_1}
\qn{3} A fountain has a fountainhead which can freely rotate on a fixed pipe without friction and leakage as shown. The flow rate in vertical pipe is Q. The cross section of outlets is S and the density of fluid is $\rho$. Find the angular velocity of the fountainhead sufficiently long time after switching on the flow. Take Q = 0.01 m$^3$/s, S = 10 cm$^2$, L = 1 m, $\rho$ = 1000 kg/m$^3$.\par
\qfig{p2_phy_s3_q03_1}
\qn{4} A disk-shaped piece of material is placed in a homogeneous magnetic field generated inside a solenoid so that the axis of the disk is parallel to the axis of the solenoid. The solenoid is supplied by an AC source with frequency w such that B(t) = B$_0$cos(wt). The disk has radius R and thickness d. Determine the resistivity of the disk (in $\Omega$ m) if the disk dissipates power P. Take R = 1 m, d = 1mm, P = 0.01 W, B$_0$ = 0.1 T, w = 100 rad/s.
\qn{5} Cart of height H = 30 cm and length $\ell$ = 40 cm must pass under the table on the horizontal floor, moving uniformly and rectilinearly (figure). A light spring of stiffness k = 50 N / m was attached to the bottom of a table top. A small load of mass m = 0.4 kg was attached to the spring. With an undeformed spring, the load was at a height h = 42 cm above the floor. Then the load was released. With which minimum speed (in m/s) cart can move so that having passed under the table, it won't touch the load?\par
\qfig{p2_phy_s3_q05_1}
\qn{6} A 2.00 m long stick is flexible and is not hinged or clamped at any point. If it is made to oscillate in its second overtone, what is its wavelength?
\qn{7} A transverse wave is passing through a string shown in figure. Mass density of the string is 1 kg/m$^3$ and cross section area of string is 0.01 m$^2$. Equation of wave in string is y = 2 sin (20t $-$ 10x). Find the hanging mass (in gm), (t is time in sec and x is position in m)\par
\qfig{p2_phy_s3_q07_1}
\qn{8} Sound from two coherent sources S$_1$ and S$_2$ are produced in phase, and detected at (point P, equidistant from both the sources. Speed of sound in normal air is 340 m/sec, but in some part in path s, there is a zone of Hot air having temperature 4 times the normal temperature, and width (d = 85 m). What should be minimum frequency of sound, so that minimum intensity can be found at P? (in Hz)\par
\qfig{p2_phy_s3_q08_1}
\partsec{2}{CHEMISTRY}
\sect{SECTION-I(i)}
\qn{1} Identify complex(es) species which have total three stereoisomers and all stereoisomers have plane of symmetry.
\begin{qoptions}
\opt{A}{[Pt(NH$_3$)(py)(Br)(Cl)]}
\opt{B}{[PtCl$_2$(en)$_2$]$^{+2}$}
\opt{C}{[Pd(NH$_3$)$_3$(H$_2$O)$_2$Br]$^{+3}$}
\opt{D}{[Co(H$_2$O)$_3$F$_3$]}
\end{qoptions}
\qn{2} Consider free CO ; [Ni(CO)$_4$] ; [Co(CO)$_4$]$^-$ and [Fe(CO)$_4$]$^{2-}$ and select \textbf{CORRECT} option(s).
\begin{qoptions}
\opt{A}{Free CO have smallest bond length}
\opt{B}{[Fe(CO)$_4$]$^{2-}$ have lowest streching frequency $\upsilon$ for C--O bond}
\opt{C}{All three complexes follow Sidwic effective atomic number rule}
\opt{D}{[Fe(CO)$_4$]$^{2-}$ have strongest M(d$\pi$)--CO($\pi^*$) attachment}
\end{qoptions}
\qfullbegin
\qn{3} Among the following resonating structures, in which option(s), (I) is more stable than (II)\par
\qfig{p2_chem_s1_q03_B}
\qn{4} Reaction(s) in which inversion of configuration takes place is/are?\par
\qfig{p2_chem_s1_q04_A}
\qfullend
\qn{5} Liquids `P' and `Q' form a solution resulting minimum boiling azeotrope. Which of the following is/are correct parameters for the solution?
\begin{qoptions}
\opt{A}{$\Delta H_{mix}$ = +ve}
\opt{B}{$\Delta S_{mix}$ = +ve}
\opt{C}{$|\Delta S_{surr}|<|\Delta S_{mix}|$}
\opt{D}{$\Delta V_{mix}$ = $-$ve}
\end{qoptions}
\qn{6} Select the correct option(s):
\begin{qoptions}
\opt{A}{When intermolecular attraction dominates over repulsive forces, the gas becomes more compressive than ideal gas.}
\opt{B}{When a gas is less compressible than ideal gas, Z $<$ 1.}
\opt{C}{An ideal gas obeys ideal gas equation only at low pressure and high temperature conditions.}
\opt{D}{Gases above critical temperature are called supercritical fluids.}
\end{qoptions}
\sect{SECTION-I(ii)}
\qn{7} Match correctly the nature of radioactive isotopes in List-II with their decay mode in List-I.\par
\noindent{\footnotesize\setlength{\tabcolsep}{2pt}\begin{tabular}{|L{.8cm}|L{3cm}|L{.8cm}|L{2.8cm}|}\hline
\multicolumn{2}{|c|}{List-I} & \multicolumn{2}{c|}{List-II}\\\hline
(P) & Alpha decay & (1) & Al$^{24}$\\\hline
(Q) & Electron (beta decay) & (2) & Al$^{29}$\\\hline
(R) & Positron decay & (3) & Na$^{20}$\\\hline
(S) & Gamma decay & (4) & U$^{238}$\\\hline
 & & (5) & Mg$^{26}$\\\hline
\end{tabular}}
\begin{qoptions}
\opt{A}{P $\to$ 4; Q $\to$ 2, 5; R $\to$ 1, 3 ; S $\to$ 1, 2, 3, 4, 5}
\opt{B}{P $\to$ 4; Q $\to$ 1, 3; R $\to$ 2, 5; S $\to$ 1, 2, 3, 4, 5}
\opt{C}{P $\to$ 4; Q $\to$ 2, 3; R $\to$ 1, 5; S $\to$ 1, 2, 3, 4, 5}
\opt{D}{P $\to$ 4; Q $\to$ 1, 5; R $\to$ 2, 5 ; S $\to$ 1, 2, 3, 4, 5}
\end{qoptions}
\qn{8} Consider the example of $\pi$-acid ligands given in List-II and their donor ability during complex formation in List-I find \textbf{CORRECT} code of matching :-\par
\noindent{\footnotesize\setlength{\tabcolsep}{2pt}\begin{tabular}{|L{.8cm}|L{3cm}|L{.8cm}|L{2.8cm}|}\hline
\multicolumn{2}{|c|}{\parbox{3.6cm}{\centering List-I (Type of $\pi$-acid ligand)}} & \multicolumn{2}{c|}{List-II (Example)}\\\hline
(P) & Anionic 6e$^-$ donor & (1) & Nitrosonium\\\hline
(Q) & Cationic 2e$^-$ donor & (2) & Triethylphosphine\\\hline
(R) & Neutral 2e$^-$ donor & (3) & Methylcyanide\\\hline
(S) & Anionic 2e$^-$ donor & (4) & Cyanide\\\hline
 & & (5) & Cyclopentadienyl\\\hline
 & & (6) & Carbonyl\\\hline
\end{tabular}}\par
The \textbf{CORRECT} option is
\begin{qoptions}
\opt{A}{P $\to$ 5 ; Q $\to$ 1,3 ; R $\to$ 2,6 ; S $\to$ 4}
\opt{B}{P $\to$ 5 ; Q $\to$ 1 ; R $\to$ 2,3,6 ; S $\to$ 4}
\opt{C}{P $\to$ 3,5 ; Q $\to$ 1 ; R $\to$ 2 ; S $\to$ 4,6}
\opt{D}{P $\to$ 4,6 ; Q $\to$ 1,2 ; R $\to$ 3 ; S $\to$ 5}
\end{qoptions}
\qn{9} Consider the product of reactions given in List-I and specification about product in reactions in List-I find \textbf{CORRECT} code of matching :-\par
\noindent{\footnotesize\setlength{\tabcolsep}{2pt}\begin{tabular}{|L{.8cm}|L{2.4cm}|L{.8cm}|L{3.4cm}|}\hline
\multicolumn{2}{|c|}{\parbox{3.2cm}{\centering List-I (Specification about product of reaction)}} & \multicolumn{2}{c|}{List-II (Reactions)}\\\hline
(P) & Green coloured residue or solution & (1) & $\mathrm{K_2Cr_2O_7}\xrightarrow{\mathrm{H_2SO_3}}$\\\hline
(Q) & Paramagnetic (any product) & (2) & Pyrolusite $\xrightarrow[\Delta,\ \text{fused}]{\mathrm{KOH+O_2}}$\\\hline
(R) & Tetrahedral geometry at central metal & (3) & $\mathrm{CuSO_{4(aq.)}}\xrightarrow[\mathrm{HCl}]{\text{conc.}}$\\\hline
(S) & Redox change & (4) & $\mathrm{K_2CrO_{4(s)}}\xrightarrow[\mathrm{HNO_3}]{\text{conc.}}$\\\hline
 & & (5) & $\mathrm{K_4[Fe(CN)_6]}\xrightarrow[\mathrm{H_2O_2}]{\mathrm{H_2SO_4}}$\\\hline
\end{tabular}}\par
The \textbf{CORRECT} option is
\begin{qoptions}
\opt{A}{P $\to$ 1,2,5; Q $\to$ 1,2,3,5; R $\to$ 1,2,3,4; S $\to$ 2,3,4}
\opt{B}{P $\to$ 1,2; Q $\to$ 1,2,4; R $\to$ 3,4,5; S $\to$ 2,3,5}
\opt{C}{P $\to$ 1,2,3; Q $\to$ 1,2,3,5; R $\to$ 2,3,4; S $\to$ 1,2,5}
\opt{D}{P $\to$ 2,3; Q $\to$ 1,2,3,4,5; R $\to$ 2,3,4; S $\to$ 1,2,3}
\end{qoptions}
\qn{10} Match the List I with List II and select the correct answer for major product using the codes given below the Lists.\par
\noindent{\footnotesize\setlength{\tabcolsep}{2pt}\begin{tabular}{|L{.7cm}|L{4.0cm}|L{.7cm}|L{2.2cm}|}\hline
\multicolumn{2}{|c|}{List - I} & \multicolumn{2}{c|}{List - II}\\\hline
(P) & $\mathrm{(CH_3)_3C\text{--}Br}\xrightarrow[\Delta]{\mathrm{C_2H_5\ddot{O}^{\ominus}Na^{\oplus}}}$ & (1) & Elimination Reaction\\\hline
(Q) & $\mathrm{(CH_3)_3C\text{--}OH}\xrightarrow{\mathrm{H^+/\Delta}}$ & (2) & Carbocation intermediate\\\hline
(R) & $\mathrm{(CH_3)_3C\text{--}Br}\xrightarrow{\mathrm{H_2O}}$ & (3) & Reaction follows I$^{\text{st}}$ order kinetics\\\hline
(S) & $\mathrm{CH_3\text{--}Br}\xrightarrow{\mathrm{CH_3CH_2\ddot{O}^{\ominus}Na^{+}}}$ & (4) & Reaction follows II$^{\text{nd}}$ order kinetics\\\hline
\end{tabular}}\par
The \textbf{CORRECT} option is
\begin{qoptions}
\opt{A}{P $\to$ 1, 2, 4 ; Q $\to$ 1, 2,3 ; R $\to$ 2,3 ; S $\to$ 4}
\opt{B}{P $\to$ 1,2, 4 ; Q $\to$ 1, 2,3 ; R $\to$ 2,3 ; S $\to$ 3,4}
\opt{C}{P $\to$ 1,2 ; Q $\to$ 1, 3 ; R $\to$ 2,3 ; S $\to$ 4}
\opt{D}{P $\to$ 1, 4 ; Q $\to$ 1, 2,3 ; R $\to$ 2,3 ; S $\to$ 4}
\end{qoptions}
\sect{SECTION-II}
\qn{1} 2-Butanol $\xrightarrow[\Delta]{\mathrm{H^+}}$ Products\par
If all products formed in the above reaction undergo reaction with Br$_2$(CCl$_4$) individually then number of products formed will be :
\qfullbegin
\qn{2} How many reactions among the following are giving alcohol as one of the product through S$_N$1 mechanism.\par
\qfig{p2_chem_s3_q02_1}
\qn{3} How many reactions giving 2$^\circ$ alcohol as one of the product.\par
\qfig{p2_chem_s3_q03_1}
\qfullend
\qn{4} For a first order gaseous reaction occurring at constant volume and temperature aA $\to$ (b - 1) B + C (with b $>$ a), the pressure of the system rose by $50\left(\dfrac ba-1\right)$ \% in 10 min. The half life of reaction (in min) is.
\qn{5} In the estimation of nitrogen content in the amino acid, valine, 50 ml of 0.15M-HClO$_4$ solution is added in a sample of valine, in glacial acetic acid.\par
$\mathrm{(CH_3)_2C(NH_2)COOH}$ + $\mathrm{HClO_4}$ $\to$\par $\mathrm{(CH_3)_2C(NH_3^+)COOH}$ + $\mathrm{ClO_4^-}$\par
The unreacted perchloric acid requires 20 ml of 0.175M solution of sodium acetate. Calculate the volume of N$_2$(g) released (in ml) at STP, when the same quantity of valine is reacted with nitrous acid?\par
$\mathrm{(CH_3)_2C(NH_2)COOH}$ + $\mathrm{HNO_2}$ $\to$ $\mathrm{(CH_3)_2C(OH)COOH}$ + $\mathrm{N_2}$ + $\mathrm{H_2O}$
\qn{6} At 25$^\circ$C, the vapour pressure of water increases from 30 mm Hg to 31 mm Hg on increasing the external pressure up to 10 atm. If the vapour pressure of water is to be maintained at 30 mm Hg at 25$^\circ$C and 10 atm external pressure, how many gram NaCl should be added per kg of water?
\qn{7} Number of species that do not exist or are unstable among the following.\par
$\mathrm{VOF_3}$, $\mathrm{CuI_2}$, $\mathrm{MnF_7}$, $\mathrm{Cr_2O_3}$, $\mathrm{Cu_2I_2}$, $\mathrm{ZnCl_2}$, $\mathrm{FeI_3}$, $\mathrm{FeBr_2}$, $\mathrm{MnO_2}$, $\mathrm{CrO(O_2)_2}$, $\mathrm{Sc^{+4}}$
\qfullbegin
\qn{8} As per CFT consider the following different spin state according to ligand field strength. (In octahedral complexes)\par
\qfig{p2_chem_s3_q08_1}
\qfullsplit
(P) = Number of spin state which is/are can be obtained by irrespective nature of ligand.\par
(Q) = Number of spin state which can be obtained only when CFSE $>$ PE but not in other cases.\par
\textbf{Find sum of P \& Q.}
\qfullend
\partsec{3}{MATHEMATICS}
\sect{SECTION-I(i)}
\qn{1} 8 players compete in a tournament and there is exactly one match between every two players. The winner gets 1 point, loser 0 point and each gets $\dfrac12$ point if the game is drawn. The final result is known where everyone gets different score and player placed second gets the same score as the total of four bottom players, then which of the following statements is/are correct ?
\begin{qoptions}
\opt{A}{The total of the player placing second is 6}
\opt{B}{The total score of the top four players is 22}
\opt{C}{The score of winner is 7}
\opt{D}{The total score of third \& fourth players is 9.}
\end{qoptions}
\qn{2} If N$_1$ denotes the number of ways in which 6 different prizes can be distributed among 4 players if each player recieve any number of prizes and N$_2$ denotes the number of ways when each player recieves atleast one prize then
\begin{qoptions}
\opt{A}{N$_1 = 6^4$}
\opt{B}{N$_1 = 4^6$}
\opt{C}{N$_2 = 1560$}
\opt{D}{N$_2 = 2536$}
\end{qoptions}
\qn{3} $ax^2+bx+a=0$ ($a,b\in R$) has roots $\alpha,\beta$. Consider the family of lines $(b^2-2a^2+c^2)x+8a^2=2(b^2-c^2)y$ ($c$ real). If $|\alpha-\beta|=2$ then
\begin{qoptions}
\opt{A}{fixed point $\left(-\dfrac47,\dfrac27\right)$}
\opt{B}{fixed point $(2,-4)$}
\opt{C}{number of lines of the family making a triangle of area 16 with the coordinate axes is only 4}
\opt{D}{number of lines of the family making a triangle of area 2 is only 2}
\end{qoptions}
\qn{4} Pair of tangents from P$(5,-1)$ to $x^2+y^2-2x-6y-6=0$ meeting at A and B. O is the centre.
\begin{qoptions}
\opt{A}{Length of AB $=4\sqrt2$ units}
\opt{B}{Area of $\Delta$ABP $=8$ square unit}
\opt{C}{Equation of circumcircle of $\Delta$OAP is $x^2+y^2-6x-2y+2=0$}
\opt{D}{$\angle$PAB $=45^\circ$}
\end{qoptions}
\qn{5} S: $x^2+y^2-2ax-2by=0$ is a given circle and $n$ is the number of chords of the circle drawn from $(2a,0)$ which are bisected by the y-axis.
\begin{qoptions}
\opt{A}{If $b^2=16a^2$, $n=3$}
\opt{B}{If $b^2=12a^2$, $n=2$}
\opt{C}{If $b^2=8a^2$, $n=1$}
\opt{D}{If $b^2=4a^2$, $n=0$}
\end{qoptions}
\qn{6} $\alpha x+y+z=0$, $x-\beta y+3z=0$ and $x-y-z=0$ has a non trivial solution, then
\begin{qoptions}
\opt{A}{$\alpha=1$ or $\beta=-5$}
\opt{B}{$\alpha=-1$ or $\beta=3$}
\opt{C}{$\alpha=-1$ or $\beta=-3$}
\opt{D}{$\alpha=1$ or $\beta=3$}
\end{qoptions}
\sect{SECTION-I(ii)}
\qn{7} Consider the lines $\ell_1: 2x-y+3=0$, $\ell_2: y=0$ and $\ell_3: y-kx-3=0$, then the number of circles which touches all the 3 given lines, if. Match List-I with List-II.\par
\noindent{\footnotesize\setlength{\tabcolsep}{2pt}\begin{tabular}{|L{.8cm}|L{3.4cm}|L{.8cm}|L{2.4cm}|}\hline
\multicolumn{2}{|c|}{\parbox{4.2cm}{\centering List-I}} & \multicolumn{2}{c|}{List-II}\\\hline
(P) & $k=0$ is(are) & (1) & infinite\\\hline
(Q) & $k=1$ is(are) & (2) & 1\\\hline
(R) & $k=2$ is(are) & (3) & 2\\\hline
(S) & $k=3$ is(are) & (4) & 4\\\hline
\end{tabular}}\par
\begin{qoptions}
\opt{A}{P$\to$2; Q$\to$3; R$\to$1; S$\to$4}
\opt{B}{P$\to$3; Q$\to$4; R$\to$1; S$\to$4}
\opt{C}{P$\to$2; Q$\to$1; R$\to$3; S$\to$4}
\opt{D}{P$\to$3; Q$\to$1; R$\to$4; S$\to$3}
\end{qoptions}
\qn{8} Match List-I with List-II.\par
\noindent{\footnotesize\setlength{\tabcolsep}{2pt}\begin{tabular}{|L{.7cm}|L{5.4cm}|L{.7cm}|L{.8cm}|}\hline
\multicolumn{2}{|c|}{List-I} & \multicolumn{2}{c|}{List-II}\\\hline
(P) & Number of points in the plane of three different coplanar non-concurrent lines which are equidistant from them, may be & (1) & 1\\\hline
(Q) & Let N denote the number of even numbers of distinct digits between 200 and 3000 that can be formed with the digits 0, 1, 2, 3, 4, then $\left\lfloor\dfrac N{10}\right\rfloor$ is equal to (where [.] denotes greatest integer function) & (2) & 2\\\hline
(R) & For the points A(0,0), B(1,0), C(0,2), D(3,4), there exist a point P such that PA+PB+PC+PD is minimum. Then number of such point P may be & (3) & 4\\\hline
(S) & Let A, B, C, D be four non-coplanar points. Then the number of planes which are equidistant from all the four points can be & (4) & 7\\\hline
 & & (5) & 5\\\hline
\end{tabular}}\par
\begin{qoptions}
\opt{A}{P$\to$3,4; Q$\to$2; R$\to$1; S$\to$5}
\opt{B}{P$\to$2; Q$\to$4,5; R$\to$1; S$\to$3}
\opt{C}{P$\to$2,3; Q$\to$5; R$\to$1; S$\to$4}
\opt{D}{P$\to$5; Q$\to$1,2; R$\to$4; S$\to$2}
\end{qoptions}
\qn{9} Let coordinates of A is $(1,5)$ point B lies on $y-x-2=0$ and point C lies on $y-2x-8=0$. Given, AB+BC+CA is minimum. Match List-I with List-II.\par
\noindent{\footnotesize\setlength{\tabcolsep}{2pt}\begin{tabular}{|L{.7cm}|L{3.6cm}|L{.7cm}|L{2.4cm}|}\hline
\multicolumn{2}{|c|}{List-I} & \multicolumn{2}{c|}{List-II}\\\hline
(P) & AB+BC+CA $=$ & (1) & $2\sqrt{13}$\\\hline
(Q) & If co-ordinate of C is $(x_1,y_1)$ then $x_1+\dfrac{25}{8}$ is equal to & (2) & $\sqrt2+\sqrt5+\sqrt{32}$\\\hline
(R) & If co-ordinate of B is $(x_2,y_2)$ then $x_2$ is equal to & (3) & $\dfrac95$\\\hline
(S) & If co-ordinate of C is $(x_1,y_1)$ then $|2x_1|$ is equal to & (4) & 2\\\hline
 & & (5) & 5\\\hline
 & & (6) & $\dfrac94$\\\hline
\end{tabular}}\par
\begin{qoptions}
\opt{A}{P$\to$2; Q$\to$5; R$\to$5; S$\to$3}
\opt{B}{P$\to$2; Q$\to$5; R$\to$4; S$\to$3}
\opt{C}{P$\to$5; Q$\to$4; R$\to$4; S$\to$5}
\opt{D}{P$\to$1; Q$\to$4; R$\to$3; S$\to$6}
\end{qoptions}
\qn{10} If $(x-1)^2+(y-4)^2=4$, then\par
\noindent{\footnotesize\setlength{\tabcolsep}{2pt}\begin{tabular}{|L{.7cm}|L{2.2cm}|L{.7cm}|L{4.2cm}|}\hline
\multicolumn{2}{|c|}{List-I} & \multicolumn{2}{c|}{List-II}\\\hline
(P) & Range of values of $(x+y)$ is & (1) & $\left[\dfrac{-4-2\sqrt{13}}3,\dfrac{-4+2\sqrt{13}}3\right]$\\\hline
(Q) & Range of values of $\dfrac yx$ is & (2) & $[2,4-\sqrt3]\cup[4+\sqrt3,6]$\\\hline
(R) & If $x\in[0,2]$, then range of values of $y$ is & (3) & $R$\\\hline
(S) & Range of values of $\dfrac{y-3}x$ is & (4) & $\left(-\infty,\dfrac{-4-2\sqrt{13}}3\right]\cup\left[\dfrac{-4+2\sqrt{13}}3,\infty\right)$\\\hline
 & & (5) & $[5-2\sqrt2,5+2\sqrt2]$\\\hline
\end{tabular}}\par
\begin{qoptions}
\opt{A}{P$\to$5; Q$\to$4; R$\to$2; S$\to$3}
\opt{B}{P$\to$4; Q$\to$5; R$\to$3; S$\to$2}
\opt{C}{P$\to$1; Q$\to$3; R$\to$5; S$\to$2}
\opt{D}{P$\to$5; Q$\to$4; R$\to$2; S$\to$1}
\end{qoptions}
\sect{SECTION-II}
\qn{1} If $\Delta=\begin{vmatrix}1&3\cos\theta&1\\ \sin\theta&1&3\cos\theta\\ 1&\sin\theta&1\end{vmatrix}$, then value of $\dfrac{\Delta_{\max}}{2}$ is.
\qn{2} Number of shortest paths from hostel to allen is equal to (as shown in given figure).\par
\qfig[0.75]{p2_math_s3_q02_1}
\qn{3} Two circles whose equations are $x^2+y^2=9$ and $x^2+y^2-8x-6y+n^2=0$, $n\in I$, have exactly two common tangents. Find the number of values of `n'.
\qn{4} Consider a family of lines $(4b+3)x-(b+1)y=2b+1$, where $b\in R$. The minimum area of the triangle which a member of this family with negative gradient makes with the positive semi-axes is equal to.
\qn{5} The sum of all the roots of equation $\begin{vmatrix}x(x-2)&1&2\\ 3&(x+2)(x-4)&4\\ 5&6&(x+4)(x-6)\end{vmatrix}=0$ will be-.
\qn{6} Images of circle C: $x^2+y^2-8x-10y+16=0$ are taken in variable line $(\lambda+1)x+(1-2\lambda)y+\lambda=2$ $\forall\lambda\in R$. If locus of centre of there images intersects circle `C' at A and B, then $\dfrac{AB}{\sqrt3}$, is.
\qn{7} Let M be any point M on the arc AB of circumcircle of right angled triangle ABC and the maximum value of AM+BM+CM is $\sqrt{29N}$, where $a=3$, $b=4$, $c=5$. Then N is $\ldots$
\qn{8} A palindrome is a whole number that reads the same forwards and backwards. If one neglect the colon, certain times displayed on a digital watch (12-hour clock) palindromes. Examples are 2:02, 3:33 and 10:01. If N is number of times during a 12-hour period are palindromes, then $\dfrac N{10}$ is equal to.
\end{multicols}
\end{document}
